% Auriga theme
% https://github.com/anishathalye/auriga

\documentclass[14pt,aspectratio=169]{beamer}
\usepackage{pgfpages}
\usepackage{fancyvrb}
\usepackage{tikz}
\usepackage{pgfplots}

\newif\ifnotes % \notestrue to show speaker notes on a second screen
\ifnotes
\setbeamertemplate{note page}[plain]
\setbeameroption{show notes on second screen=right}
\fi

\usetheme{auriga}
\usecolortheme{auriga}

% define some colors for a consistent theme across slides
\definecolor{red}{RGB}{181, 23, 0}
\definecolor{blue}{RGB}{0, 118, 186}
\definecolor{gray}{RGB}{146, 146, 146}
\definecolor{title_color}{RGB}{0, 118, 186}


\settitlecolor{title_color}
\title{Results on PBM-driven ODE discovery and ODE parameter estimation}

\author{Urh Primožič, IJS}



\begin{document}

{
  % rather than use the frame options [noframenumbering,plain], we make the
  % color match, so that the indicated page numbers match PDF page numbers
  \setbeamercolor{page number in head/foot}{fg=white}
  \begin{frame}
    \titlepage
  \end{frame}
}

\begin{frame}{ODE discovery with Probmot}
TODO delovanje in potrebe po novem sistemu 
\end{frame}
 \begin{frame}{Structure changes}
    Kar je novega v pybm, primerjava s probmotom  
 \end{frame}
 \begin{frame}{Other possible solutions}
    Vsi speedupi, ki niso implementirani, minus estimacija
 \end{frame}
 \begin{frame}{Parameter estimation}
    problem , unobserved vs observed
 \end{frame}
\begin{frame}{Integration using multishooting}
    
\end{frame}
\begin{frame}{Matching}
    
\end{frame}
\begin{frame}{Integral matching }
    
\end{frame}
\begin{frame}{Batteling caotic loss surface}
    Integral lenght, homotopic methods
\end{frame}
\begin{frame}{Collocation}
    matching on unobserved
\end{frame}
\begin{frame}{Note on MAGI}
    TODO razlike 
\end{frame}
\begin{frame}{Question: lenght of integration interval}
    
\end{frame}
\begin{frame}{Results}
    
\end{frame}
\begin{frame}{Back to ProBmot}
    
\end{frame}



\end{document}
